% =====================================================================
% PROJETO DE ENGENHARIA DE SOFTWARE — AULA 1 (26/10)
% Entender o problema e definir o escopo
% Caso 15 · Reserva de laboratórios da faculdade
% =====================================================================

\documentclass[a4paper,12pt]{article}

\usepackage[utf8]{inputenc}
\usepackage[T1]{fontenc}
\usepackage[brazilian]{babel}
\usepackage{times}
\usepackage[a4paper,top=3cm,left=3cm,right=3cm,bottom=2.5cm]{geometry}
\usepackage{xltabular}
\usepackage{array}
\usepackage{amssymb}
\usepackage{titlesec}
\usepackage{enumitem}

\setlength{\parindent}{0pt}
\setlength{\parskip}{6pt}
\renewcommand{\arraystretch}{1.3}
\setlist{itemsep=0pt,topsep=2pt}

\titleformat{\section}{\fontsize{13}{15}\bfseries}{\thesection.}{1em}{}
\titleformat{\subsection}{\fontsize{12}{14}\bfseries}{\thesubsection.}{1em}{}

\newcolumntype{Y}{>{\raggedright\arraybackslash}X}

\begin{document}
\pagestyle{plain}

\begin{center}
  {\small\bfseries PROJETO DE ENGENHARIA DE SOFTWARE}\\[4pt]
  {\fontsize{16}{20}\selectfont\bfseries Aula 1 · 26/10 · Entender o problema e definir o escopo}\\[6pt]
  {\large Caso 15 · Reserva de laboratórios da faculdade}
\end{center}

\vspace{6pt}

\begin{xltabular}{\linewidth}{|l|Y|}
  \hline
  \textbf{Professor} & Sérgio Santana \\ \hline
  \textbf{Disciplina} & Laboratório de Engenharia de Software \\ \hline
  \textbf{Integrante} & Wilia Francisco da Silva \\ \hline
  \textbf{Regra para testar} & Reservas de 1 a 4 horas, para até 40 alunos \\ \hline
\end{xltabular}

% ---------------------------------------------------------------------
\section{O problema}

A coordenação da faculdade recebia pedidos de laboratório por e-mail e
dois professores chegavam para o mesmo horário. A faculdade tem quatro
laboratórios: o Lab 202 tem 30 lugares e os demais (Lab 101, Lab 102 e
Lab 201) têm 40. A faculdade abre de segunda a sexta das 07:00 às 22:00
e aos sábados até 12:00. O sistema deve mostrar a agenda dos
laboratórios, permitir reservar e impedir choques de horário.

A equipe usou a versão 1.0 do sistema, no Laboratório de Sistemas, por
20 minutos como se fosse o cliente, e separou os requisitos que o
sistema já atende dos que ainda faltam.

% ---------------------------------------------------------------------
\section{Requisitos funcionais}

A coluna \emph{Regras} indica as regras de negócio da tela do sistema
(RN01 a RN18) ligadas a cada requisito. Os requisitos RF06 a RF10 não
existem na versão 1.0 e, por isso, ainda não têm regras na tela.

\begin{xltabular}{\linewidth}{|c|Y|>{\centering\arraybackslash}p{2.3cm}|>{\centering\arraybackslash}p{2.3cm}|}
  \hline
  \textbf{Código} & \textbf{Requisito} & \textbf{Regras} & \textbf{Versão 1.0} \\ \hline
  RF01 & Cadastrar uma reserva com professor, disciplina, laboratório, data, horário de início, duração e número de alunos. & RN01 a RN06, RN09 a RN13 & Atende \\ \hline
  RF02 & Impedir reservas com horários sobrepostos no mesmo laboratório. & RN07 & \textbf{Atende\newline parcialmente}$^{1}$ \\ \hline
  RF03 & Cancelar ou excluir uma reserva, liberando o horário para outros professores. & RN08, RN18 & \textbf{Atende\newline parcialmente}$^{2}$ \\ \hline
  RF04 & Listar as reservas por data e horário e buscar pelo nome do professor ou pela disciplina. & RN14, RN15, RN17 & \textbf{Atende\newline parcialmente}$^{3}$ \\ \hline
  RF05 & Mostrar a ocupação de cada laboratório no dia seguinte. & RN16 & \textbf{Atende\newline parcialmente}$^{4}$ \\ \hline
  RF06 & Exibir a agenda de cada laboratório em grade (dia $\times$ horário), destacando os horários livres. & -- & \textbf{Falta} \\ \hline
  RF07 & Consultar a disponibilidade de um laboratório antes de preencher a reserva. & -- & \textbf{Falta} \\ \hline
  RF08 & Editar uma reserva já feita (data, horário ou laboratório). & -- & \textbf{Falta} \\ \hline
  RF09 & Identificar o usuário por login, separando o perfil de professor do perfil da coordenação. & -- & \textbf{Falta} \\ \hline
  RF10 & Notificar o professor por e-mail quando a reserva for confirmada ou cancelada. & -- & \textbf{Falta} \\ \hline
\end{xltabular}

{\small
$^{1}$ RF02 (RN07): bloqueia apenas reservas com o mesmo horário de início.
Reservas cujos horários se cruzam são aceitas (exemplo: Lab 101, 06/11/2026,
das 19:00 às 22:00 e das 20:00 às 21:00).

$^{2}$ RF03 (RN08): a reserva é cancelada, mas o horário não é liberado. Uma
nova reserva no mesmo laboratório, data e horário é recusada com a mensagem
``já está reservado nesse horário''.

$^{3}$ RF04 (RN15): a busca encontra pelo nome do professor, mas não pela
disciplina.

$^{4}$ RF05 (RN16): o painel Ocupação de amanhã soma também as horas das
reservas canceladas.
}

% ---------------------------------------------------------------------
\section{Requisitos não funcionais (qualidades)}

\begin{xltabular}{\linewidth}{|c|l|Y|}
  \hline
  \textbf{Código} & \textbf{Qualidade} & \textbf{Descrição} \\ \hline
  RNF01 & Usabilidade & Ser fácil de usar no celular: um professor faz uma reserva em menos de 2 minutos, sem treinamento. \\ \hline
  RNF02 & Desempenho  & Mostrar a agenda e o resultado da busca em até 2 segundos. \\ \hline
  RNF03 & Segurança   & Permitir que apenas a coordenação exclua ou altere reservas de outros professores. \\ \hline
\end{xltabular}

% ---------------------------------------------------------------------
\section{O que fica fora do sistema}

\begin{itemize}[leftmargin=1.2cm]
  \item Controle de presença dos alunos nas aulas.
  \item Manutenção e inventário dos computadores dos laboratórios.
  \item Reserva de outros espaços da faculdade, como auditório e salas comuns.
  \item Integração com o sistema acadêmico da faculdade (matrículas e notas).
  \item Acesso de alunos ao sistema (consulta ou reserva de laboratórios).
\end{itemize}

% ---------------------------------------------------------------------
\section{Histórias de usuário}

Cada história foi escrita a partir de um dos requisitos da Seção 2.

\begin{xltabular}{\linewidth}{|c|Y|c|}
  \hline
  \textbf{Nº} & \textbf{História} & \textbf{Requisito} \\ \hline
  H01 & Como \textbf{professor}, quero ver a agenda de um laboratório, para escolher um horário livre sem tentar várias vezes. & RF06 \\ \hline
  H02 & Como \textbf{professor}, quero reservar um laboratório pelo celular, para não precisar mandar e-mail para a coordenação. & RF01 \\ \hline
  H03 & Como \textbf{professor}, quero cancelar minha reserva, para liberar o laboratório para outro colega. & RF03 \\ \hline
  H04 & Como \textbf{professor}, quero alterar a data ou o horário da reserva, para não precisar cancelar e criar outra. & RF08 \\ \hline
  H05 & Como \textbf{coordenação}, quero buscar as reservas de um professor ou disciplina, para resolver dúvidas rapidamente. & RF04 \\ \hline
  H06 & Como \textbf{professor}, quero receber um aviso quando a reserva for confirmada, para ter certeza de que o laboratório está garantido. & RF10 \\ \hline
\end{xltabular}

\end{document}
